% Self-contained section fragment. Requires amsmath, amssymb, amsthm.
% Suggested source citations: DLMF 25.12.12 and the repository's
% Polylogarithm_Resonance_And_Inverse_Transseries, prop:noninteger.
\section{Exact ramification of the inverse at the polylogarithm endpoint}
\label{ep:sec}

The endpoint of a forward function need not be a singular point of its
inverse. For the polylogarithm this distinction produces an infinite
discrete family of removable endpoints. We give an exact classification,
including the smallest branched cover that resolves every noninteger rational order.
The convergent nonintegral charts underlying the result already occur in
Proposition~\texttt{prop:noninteger} of the repository report~\cite{RepoInverse}.
Our contribution here is the minimal ramification degree, the resulting
if-and-only-if removability criterion, and explicit Taylor data at the
exceptional orders. The general use of Puiseux uniformization is classical.

Fix a real number $s>1$, put $\rho_s=\zeta(s)$, and let $G_s$ denote
the inverse germ of $\Li_s$ continued from the interval
$0<w<\rho_s$. Since $\Li_s$ increases strictly from $0$
to $\rho_s$ on $0<z<1$, this specifies the branch without a complex
sheet convention. Near its endpoint write
\[
 y=\rho_s-w,\qquad t=-\log G_s(\rho_s-y),\qquad
 \Delta_s(t)=\rho_s-\Li_s(e^{-t}).
\]
Initially $y,t>0$. All subsequent identities refer to the analytic
continuations of this local branch. In particular, a continuation of
$G_s$ through $\rho_s$ need not remain an inverse of the principal
polylogarithm on both sides of its chosen cut.

We say that the endpoint has finite ramification degree $d$ if $d$ is
the smallest positive integer for which
$G_s(\rho_s-\eta^d)$ extends holomorphically to $\eta=0$.
The substitution is made on the branch agreeing with positive $\eta$.
Degree one means that the endpoint is removable for the inverse.
When no such integer exists, we say that the endpoint has infinite
ramification. Equivalently, the bounded local inverse is not resolved by
any finite branched cover; thus it cannot have finitely many local
monodromy branches around the endpoint.

For nonintegral $s$, the classical convergent Jonqui\`ere expansion~\cite{DLMF} gives
\begin{equation}
 \Delta_s(t)=c_s t^{\alpha}+A_s(t),\qquad
 \alpha=s-1,\quad c_s=-\Gamma(-\alpha),
 \label{ep:eq:normal}
\end{equation}
where
\begin{equation}
 A_s(t)=\sum_{k\ge1}a_k(s)t^k,
 \qquad a_k(s)=\frac{(-1)^{k+1}}{k!}\zeta(s-k),
 \qquad |t|<2\pi.
 \label{ep:eq:analytic}
\end{equation}
The power in \eqref{ep:eq:normal} is taken on the selected logarithmic
cover. Notice that $c_s>0$ for $1<s<2$, whereas
$a_1(s)=\zeta(s-1)>0$ for $s>2$.

\begin{theorem}[Minimal endpoint ramification]
\label{ep:thm:classification}
Let $s>1$ be real.
\begin{enumerate}
\item If $s$ is nonintegral and $s-1=p/q$ in lowest positive terms,
then the endpoint ramification degree of $G_s$ is
\[
 d_s=\min\{p,q\}.
\]
\item If $s-1$ is irrational, the endpoint has infinite ramification.
\item If $s$ is an integer with $s\ge2$, the endpoint has infinite
ramification.
\end{enumerate}
Consequently, $G_s$ extends holomorphically through $w=\zeta(s)$
if and only if
\begin{equation}
 s=1+\frac1q,\qquad q=2,3,4,\ldots.
 \label{ep:eq:exceptional}
\end{equation}
\end{theorem}

\begin{proof}
The local changes $z=e^{-t}$ and $t=-\log z$ are biholomorphic near
$t=0$ and $z=1$. Thus $G_s$ and $t(y)$ have exactly the same
ramification degree.

First suppose $\alpha=p/q$ is rational and nonintegral. Put $t=v^q$.
If $p<q$, equation~\eqref{ep:eq:normal} becomes
\[
 y=v^p B(v),\qquad
 B(v)=c_s+\sum_{k\ge1}a_k(s)v^{qk-p},\qquad B(0)=c_s\ne0.
\]
Choose the analytic $p$th root of $B$ which is positive at the origin.
The map $v\mapsto\eta=vB(v)^{1/p}$ has nonzero derivative there.
Its holomorphic inverse $v=V(\eta)$ gives
\[
 y=\eta^p,\qquad t(y)=V(\eta)^q,\qquad
 G_s(\rho_s-\eta^p)=e^{-V(\eta)^q}.
\]
This proves that degree $p$ suffices. Conversely,
\begin{equation}
 t(y)=c_s^{-q/p}y^{q/p}(1+o(1)).
 \label{ep:eq:low-leading}
\end{equation}
If $t(\eta^d)$ were holomorphic at zero, its order of vanishing
would be $dq/p$, an integer. Coprimality forces $p\mid d$.
Thus $p$ is minimal.

If $p>q$, write instead
\[
 y=v^q B(v),\qquad
 B(v)=a_1(s)+c_s v^{p-q}
            +\sum_{k\ge2}a_k(s)v^{q(k-1)},\qquad B(0)>0.
\]
The map $\eta=vB(v)^{1/q}$ is again a local analytic coordinate,
so that $y=\eta^q$ and $t=V(\eta)^q$. Hence degree $q$ suffices.
To prove minimality, let $H_s$ be the analytic inverse of $A_s$ at
zero and put $a=a_1(s)$. Substitution into
\eqref{ep:eq:normal}, followed by one Taylor expansion of $A_s$,
gives
\begin{equation}
 t(y)-H_s(y)
   =-c_s a^{-\alpha-1}y^{\alpha}(1+o(1)),
 \qquad y\downarrow0.
 \label{ep:eq:high-first}
\end{equation}
Indeed, $t\sim y/a$, the difference is $O(y^{\alpha})$, and the
linear term of $A_s(t)-A_s(H_s(y))$ is
$a(t-H_s(y))(1+o(1))$. Since $c_s\ne0$, a holomorphic
$t(\eta^d)-H_s(\eta^d)$ would have vanishing order $d\alpha=dp/q$.
This forces $q\mid d$ and establishes minimality.

For irrational $0<\alpha<1$, the same leading law
$t(y)=c_s^{-1/\alpha}y^{1/\alpha}(1+o(1))$ rules out every finite
degree $d$, because $d/\alpha$ is not an integer. For irrational
$\alpha>1$, equation~\eqref{ep:eq:high-first} rules out every
finite degree because $d\alpha$ is not an integer.

It remains to check integer orders. For $s=m\ge3$, the classical
resonant expansion has the form
\begin{equation}
 \Delta_m(t)=A_m(t)+b_m t^{m-1}\log t,
 \qquad A_m'(0)=\zeta(m-1)>0,
 \qquad b_m=\frac{(-1)^{m-1}}{(m-1)!}\ne0,
 \label{ep:eq:integer}
\end{equation}
with $A_m$ analytic. If $H_m=A_m^{-1}$ and $a=\zeta(m-1)$,
the preceding first-correction argument now yields
\begin{equation}
 t(y)-H_m(y)
   =-b_m a^{-m}y^{m-1}\log y\,(1+o(1)).
 \label{ep:eq:integer-first}
\end{equation}
After $y=\eta^d$, its nonzero leading term is a constant times
$\eta^{d(m-1)}\log\eta$. No holomorphic function has this
leading behavior: its ratio to its first nonzero integral power is
a finite nonzero constant. Therefore no finite degree suffices.

At $m=2$,
\[
 \Delta_2(t)=t\log(1/t)+t+O(t^2),\qquad
 t(y)\sim\frac{y}{\log(1/y)}.
\]
If $t(\eta^d)$ were holomorphic and nonzero, its leading monomial
would be $C\eta^j$ for some integer $j\ge1$ and $C\ne0$.
The displayed asymptotic would force $j=d$, but then the leading
coefficient would be zero, a contradiction. This treats all cases.
Since a nonintegral rational $\alpha>1$ has $q\ge2$, while
$0<\alpha<1$ gives degree one precisely when $p=1$, the final
criterion follows.
\end{proof}

\begin{remark}[Scope of the monodromy statement]
For rational orders the proof explicitly constructs a holomorphic
function on a finite cover, and its minimal degree is the exact number
of local branches. In the other cases it proves an obstruction to
every finite cover. If a bounded inverse had finite local monodromy,
passing to its finite covering would make it single-valued and its
bounded singularity removable, contradicting that obstruction.
This is a local assertion. It neither identifies every singularity of
the inverse continued globally nor determines the Taylor radius at
the original point $w=0$.
\end{remark}

\subsection{Taylor charts at the removable endpoints}

At the exceptional order $s_q=1+1/q$ define
\[
 \rho_q=\zeta(s_q),\qquad
 c_q=-\Gamma(-1/q)=q\Gamma(1-1/q)>0,
\]
and the convergent analytic germ
\begin{equation}
 Y_q(v)=c_qv+\sum_{k\ge1}
        \frac{(-1)^{k+1}\zeta(1+1/q-k)}{k!}\,v^{qk}.
 \label{ep:eq:Y}
\end{equation}
Let $V_q=Y_q^{-1}$ be its analytic local inverse. Then the exact
endpoint inverse is
\begin{equation}
 G_{s_q}(\rho_q-y)=\exp\bigl(-V_q(y)^q\bigr).
 \label{ep:eq:exact-chart}
\end{equation}
The series defining $Y_q$ converges whenever $|v|^q<2\pi$;
the inverse has a positive local radius by $Y_q'(0)=c_q$.
No assertion that these two radii agree is needed.

\begin{proposition}[Finite coefficient formula and arithmetic support]
\label{ep:prop:coefficients}
Put $a_k=(-1)^{k+1}\zeta(1+1/q-k)/k!$ and write
\[
 G_{s_q}(\rho_q-y)=1+\sum_{n\ge1}g_{q,n}y^n.
\]
For every $n\ge1$,
\begin{equation}
 g_{q,n}=-\frac qn[v^{n-q}]
    e^{-v^q}\left(c_q+\sum_{k\ge1}a_kv^{qk-1}\right)^{-n},
 \label{ep:eq:lagrange}
\end{equation}
where a negative coefficient index means zero. Every extraction is
finite. The support satisfies
\begin{equation}
 \{n\ge1:g_{q,n}\ne0\}
 \subseteq q+\langle q-1,q\rangle
 =\{q+j(q-1)+kq:j,k\ge0\}.
 \label{ep:eq:support}
\end{equation}
In particular,
\begin{equation}
 g_{q,q}=-c_q^{-q},\qquad
 g_{q,2q-1}=\frac{q\zeta(1/q)}{c_q^{2q}}\ne0,
 \label{ep:eq:first-two}
\end{equation}
and $g_{q,n}=0$ for $q<n<2q-1$.
\end{proposition}

\begin{proof}
Lagrange--B\"urmann applied to $e^{-V_q(y)^q}$ gives
\[
 [y^n]e^{-V_q(y)^q}
 =\frac1n[v^{n-1}](-qv^{q-1}e^{-v^q})
                  \left(\frac v{Y_q(v)}\right)^n,
\]
which is \eqref{ep:eq:lagrange}. In its last product, the available
positive exponents are $q$ and $qk-1=(q-1)+(k-1)q$.
Their sums give \eqref{ep:eq:support}. Extracting the constant
term and then the term of degree $q-1$ proves
\eqref{ep:eq:first-two}. Finally $\zeta(1/q)<0$ because the
alternating Dirichlet eta series is positive and
$1-2^{1-1/q}<0$.
\end{proof}

For $q=2$, set $a=\zeta(1/2)$. Since $c_2=2\sqrt\pi$, the first
Taylor terms are
\begin{align}
 G_{3/2}(\zeta(3/2)-y)
 ={}&1-\frac{y^2}{4\pi}+\frac{a\,y^3}{8\pi^2}\notag\\
 &+\left(\frac1{32\pi^2}-\frac{5a^2}{64\pi^3}\right)y^4
 +O(y^5).
 \label{ep:eq:q2}
\end{align}
Thus a square-root singularity of the forward function produces a
holomorphic inverse with a quadratic critical point.

For $q=3$, put $c=3\Gamma(2/3)$ and $a=\zeta(1/3)$. Then
\begin{align}
 G_{4/3}(\zeta(4/3)-y)
 ={}&1-\frac{y^3}{c^3}+
       \frac{3a\,y^5}{c^6}+\frac{y^6}{2c^6}
       -\frac{12a^2y^7}{c^9}+O(y^8).
 \label{ep:eq:q3}
\end{align}
The missing fourth-degree term is a structural gap, explained by
\eqref{ep:eq:support}. For general $q$, derivatives of orders
$1,\ldots,q-1$ vanish at the endpoint and
\[
 G_{s_q}^{(q)}(\rho_q)=(-1)^{q+1}\frac{q!}{c_q^q}\ne0.
\]

\subsection{A correction to a proposed endpoint-circle extension}

The repository's integer-order cut-modulus theorem~\cite{RepoInverse} states that
$|\Li_s(e^u+i0)|>\zeta(s)$ for integer $s\ge2$ and
$u>0$. Its further-research question~9 asks for an all-real-order
endpoint-circle theorem and points to the integer cut-modulus argument.
The corresponding all-real extension of this cut-modulus ingredient
fails below $s=3/2$. This does not contradict the proved integer
theorem or by itself settle the full endpoint-circle question.

\begin{proposition}[Sharp local threshold for the cut-modulus inequality]
\label{ep:prop:cut-counterexample}
For fixed $1<s<3/2$, one has
\[
 |\Li_s(e^u+i0)|<\zeta(s)
\]
for all sufficiently small $u>0$. For every fixed $s\ge3/2$, the
opposite strict inequality holds for all sufficiently small $u>0$.
This proposition concerns a neighborhood of the endpoint, rather
than the entire cut.
\end{proposition}

\begin{proof}
For $1<s<2$, put $\alpha=s-1\in(0,1)$ and $\rho=\zeta(s)$.
On the upper bank $t=-u-i0$, so Jonqui\`ere's formula gives
\[
 \Li_s(e^u+i0)
 =\rho+\Gamma(-\alpha)e^{-i\pi\alpha}u^\alpha
       +\zeta(\alpha)u+O(u^2).
\]
If $\alpha\ne1/2$, then
\[
 |\Li_s(e^u+i0)|^2-\rho^2
 =2\rho\Gamma(-\alpha)\cos(\pi\alpha)u^\alpha
      +o(u^\alpha).
\]
Since $\Gamma(-\alpha)<0$, its sign is negative for
$0<\alpha<1/2$ and positive for $1/2<\alpha<1$.
At $\alpha=1/2$, the leading fractional term is purely imaginary,
and the next calculation gives
\begin{equation}
 |\Li_{3/2}(e^u+i0)|^2-\zeta(3/2)^2
 =\bigl(4\pi+2\zeta(3/2)\zeta(1/2)\bigr)u+O(u^2).
 \label{ep:eq:cut-critical}
\end{equation}
The coefficient is strictly positive. For an elementary certificate,
$\zeta(3/2)<3$ follows by integral comparison, and
$-2<\zeta(1/2)<0$ follows from the fifth alternating partial sum
\[
 0<\eta(1/2)<1-2^{-1/2}+3^{-1/2}-\tfrac12+5^{-1/2}
                   <2(\sqrt2-1).
\]
The last strict inequality may be checked by elementary squaring;
for example the rational bounds
$\sqrt2>1414/1000$, $1/\sqrt2>707/1000$,
$1/\sqrt3<578/1000$, and $1/\sqrt5<448/1000$
already prove it. Consequently the coefficient in
\eqref{ep:eq:cut-critical} exceeds $4\pi-12>0$.
For $s>2$, the boundary increment is
$\zeta(s-1)u+o(u)$ with positive real coefficient, and for $s=2$
it is $u(1-\log u)+i\pi u+O(u^2)$.
These establish the remaining local cases.
\end{proof}

For example, at $s=5/4$ and $u=10^{-6}$, a 60-digit evaluation gives
$|\Li_{5/4}(e^u+i0)|-\zeta(5/4)
\approx-0.108266505478287898$. The sign is proved by the preceding
asymptotic argument; the decimal is only a diagnostic.
The modulus gap concerns the boundary of the principal-sheet image.
A regular point on an image cut can be crossed by analytic continuation
of the inverse. Therefore the counterexample alone does not locate
the true nearest singularity of the inverse germ at $w=0$.
